% Nouveau EGA V (cote 25): archive cote 25, site page 71 (PDF page 72).
% typescript. Draft compared: none.
% AI-transcribed by Claude Opus 5.5 (claude-opus-5-5, adaptive thinking, effort high), Message Batches API, batch ega25b (msgbatch_01A7n5YuXubUPNK1zsepi8RG), 2026-11-13; not yet checked by eye.
\scanpage{25/71}
\mspage{66}

\xout{Pour autant que je puisse voir, le théorème 8.7 , notamment}
\xout{exact}
\xout{8.7 \unsure{e)}, n'admet pas d'analogue évident pour m quelconque,}

Je pense que essentiellement tous les dév\unsure{el}oppements du n°8 à
partir de 8.6. peuvent s'apdapter \inserted{également} au cas des sections linéaires avec
m quelconque. Le faire en forme constituerait un exercice sans doute
assez long et fastidieux. Je dois admettre que je ne connais pas d'appli-%
cation s'appuyant de façon essentielle sur l'analyse de cette situation
plus générale, donc nous ne sommes pas vraiment obligés d'inclure ces
développements dans les Eléments. D'autre\unsure{x} part, l'expérience prouve que
le fait de rédiger dans un cadre plus général oblige parfois à mieux
dévisser et fait mieux comprendre le fourbis, et souvent sans beaucoup
plus de mal, de plus un certain nombre d'exercices de syntaxe dans un
contexte proprement géométrique comme ici ne ferait pas de mal,-et
bien entendu il n'est pas du tout exclu qu'on en ait quand même besoin
un jour et soit content de trouver la chose faite. Aussi je te laisse
l'entière décision à ce sujet, et me borne à esquisser simplement les
\xout{\illegible{1}}énoncés qu'on pourrait éventuellement donner à ce propos.

Supposons encore que $f:X \hand{$\longrightarrow$} P$ soit non ramffié et X soit lisse sur
S,\inserted[typed]{à composantes de dimension \hand{$\geq$} m.}\hand{.} \hand{(}Dans $\underline{X}^{(m)}$ \inserted{$=V^o$}\xout{xxx} on distingue alors le sous-préschéma $\underline{X}^{(m)\mathrm{sing}} = V^1$
des zéros \unsure{singuliers} de $\phi^m$ relativement à Grass$^m$ , qui est aussi formé
géométriquement des couples $(x,\hand{$\xi$})$ tels que la variété linéaire $L_{\hand{$\xi$}}$
coupe excédentairement l'espace tangent à X en x (considéré comme
sous-variété linéaire de P), i.e. tels que les deux espaces n'engendrent
pas P tout entier. Contrairement à ce qui se passe pour m=1, lorsque
m \inserted[typed]{est quelconqué}\xout{\illegible{3}}, le morphisme $\underline{X}^{(m)}$ sing $\hand{$\longrightarrow$}$ X n'est en général pas
lisse, car la variété des L qui \inserted{\unsure{passant par x et qui}}\hand{(}coupent excédentairement une sous-varié-%
té linéaire donnée T\hand{\unsure{$\ni x$}} n'est en général pas lisse sur k: cette variété

\margin{71}
